% ===================================================================
% preamble.tex - Praeambel des Vollberichts GEA Group
% Uebernommen aus Business_Analysen/Dual_Champions und auf einen
% einzelnen Titel zurueckgebaut: keine Teile, keine Paartabellen.
% ===================================================================
\usepackage[T1]{fontenc}
\usepackage[utf8]{inputenc}
\usepackage{lmodern}
\usepackage[ngerman]{babel}
\usepackage{csquotes}
% patch ohne "footnote": microtype kann seinen Fussnotenpatch nicht anlegen,
% wenn footmisc (para) den Fussnotenapparat ersetzt.
\usepackage[patch={item,toc,eqnum}]{microtype}
\usepackage{xcolor}
\usepackage{amsmath}
\usepackage{amssymb}
\usepackage{booktabs}
\usepackage{longtable}
\usepackage{array}
\usepackage{tabularx}
\usepackage{siunitx}
\usepackage{graphicx}
\usepackage{tikz}
\usepackage{placeins}
\usepackage{needspace}
\usepackage{pgfplots}
\usepackage[hidelinks]{hyperref}
\usepackage{url}

\areaset[current]{463.07pt}{636.6pt}

% Schusterjungen und Hurenkinder verbieten (im deutschen Satz Fehler).
\clubpenalty=10000
\widowpenalty=10000
\displaywidowpenalty=10000

% Fussnotenapparat flach halten: alle Fussnoten einer Seite in einem
% durchlaufenden Absatz. Das zweispaltige dblfnote laeuft mit pgfplots und
% hyperref in eine Endlosschleife der Ausgaberoutine.
\usepackage[para]{footmisc}

\usepackage[automark]{scrlayer-scrpage}
\pagestyle{scrheadings}
\renewcommand{\sectionmark}[1]{\markright{\thesection\quad #1}}
\clearpairofpagestyles
\setkomafont{pagehead}{\normalfont\small}
\ihead{\rightmark}
\ohead{Neste}
\KOMAoptions{headsepline=true}
\cfoot*{\pagemark}

\setcounter{tocdepth}{2}

\pgfplotsset{compat=1.18}
\usepgfplotslibrary{dateplot}
\pgfplotsset{/pgf/number format/use comma, /pgf/number format/1000 sep={}}

\definecolor{kursfarbe}{RGB}{25,55,120}
\definecolor{auftragsfarbe}{RGB}{0,120,60}

% --- Zahlenformat: deutsches Dezimalkomma ---------------------------
\sisetup{
  output-decimal-marker = {,},
  group-separator       = {.},
  group-minimum-digits  = 4,
  % Ohne group-digits=integer gruppiert siunitx auch die Nachkommastellen:
  % aus 43.4679 wurde im Satz "43,467.9".
  group-digits          = integer,
  detect-weight         = true,
  detect-family         = true
}

% --- Zitierweise ------------------------------------------------------
% \Q{Schluessel}{Seite}: Fussnote mit Kurztitel (Sprung ins Quellenverzeichnis)
% und GEDRUCKTER Seite. \QW{Titel}{Schluessel}{Abrufdatum}: Webquelle; die
% Adresse steht nur im Quellenverzeichnis, damit jede URL genau einmal vorkommt.
\definecolor{quellenlink}{RGB}{20,60,120}
\newcommand{\quellensprung}[2]{\hyperref[#1]{\textcolor{quellenlink}{#2}}}
\newcommand{\kurzGBXXV}{Neste, Annual Report 2025}
\newcommand{\kurzGBXXIV}{Neste, Annual Report 2024}
\newcommand{\kurzGBXXIII}{Neste, Annual Report 2023}
\newcommand{\kurzGBXXII}{Neste, Annual Report 2022}
\newcommand{\kurzGBXXI}{Neste, Annual Report 2021}
\newcommand{\kurzGBXX}{Neste, Annual Report 2020}
\newcommand{\kurzGBXIX}{Neste, Annual Report 2019}
\newcommand{\kurzFSRXVIII}{Neste, Financial Statements Release 2018}
\newcommand{\kurzHJ}{Neste, Half-year financial report January--June 2026}
\newcommand{\Q}[2]{\footnote{\quellensprung{quelle:#1}{\csname kurz#1\endcsname}, S.~#2.}}
% \QL{Schluessel}{Dokument}{Seite}: Fussnote mit Marke; \QR{Schluessel}: dieselbe Fussnote
% erneut referenziert (nur auf derselben gedruckten Seite, geprueft von
% scripts/check_footnote_groups.py). \QI: Beleg im Text, fuer Tabellenzellen, in denen
% Fussnoten in Gleitumgebungen verloren gehen.
\newcommand{\QL}[3]{\footnote{\label{fn:#1}\quellensprung{quelle:#2}{\csname kurz#2\endcsname}, S.~#3.}}
\newcommand{\QR}[1]{\textsuperscript{\ref{fn:#1}}}
\newcommand{\tagGBXXV}{GB\,2025}\newcommand{\tagGBXXIV}{GB\,2024}\newcommand{\tagGBXXIII}{GB\,2023}
\newcommand{\tagGBXXII}{GB\,2022}\newcommand{\tagGBXXI}{GB\,2021}\newcommand{\tagHJ}{HJ\,2026}
\newcommand{\QI}[2]{{\scriptsize[\csname tag#1\endcsname\ S.\,#2]}}
\newcommand{\QW}[3]{\footnote{\quellensprung{#2}{#1} (abgerufen am #3).}}

% --- Einheiten: Makros aus data/kennzahlen.tex sind bereits deutsch gesetzt
\newcommand{\Mio}[2]{#1\,Mio.\,#2}
\newcommand{\je}[2]{#1\,#2}
\newcommand{\Pz}[1]{#1\,\%}

% --- Drei Farben, drei Verantwortlichkeiten --------------------------
% Schwarz: belegte Fakten. Blau: Einschaetzung des Verfassers. Rot: Votum.
\definecolor{vdrot}{RGB}{155,25,30}
\newcommand{\vd}[1]{{\color{vdrot}#1}}
\definecolor{bkblau}{RGB}{20,60,150}
\newcommand{\bk}[1]{{\color{bkblau}#1}}

% --- Generierte Tabellenkoerper ohne angehaengtes \relax --------------
\makeatletter
\newcommand{\tabellenkoerper}[1]{\@@input #1.tex }
\makeatother

\newcommand{\annahme}[1]{%
  \par\medskip\noindent\fbox{\parbox{0.97\linewidth}{\small\textbf{Annahme:} #1}}\par\medskip}

% =====================================================================
% Tabellen. Jeder Koerper kommt aus data/ und wird von scripts/rechnung_neste.py
% erzeugt; im Dokument steht keine Zahl von Hand.
% =====================================================================
\newcommand{\ueberblicktabelle}{%
\begin{table}[htbp]\centering\small
\caption{Kennzahlen 2025 gegen 2024 (Mio.~EUR, wo nicht anders angegeben), aus der
Mehrjahresreihe (Tabelle~\ref{tab:ergebnis}). Ver\"anderung in Einheiten der Zeile
(Prozentpunkte bei Quoten).}\label{tab:ueberblick}
\begin{tabular}{@{}lrrr@{}}\toprule
Kennzahl & 2025 & 2024 & Ver\"anderung\\\midrule
\tabellenkoerper{data/ueberblick}
\bottomrule\end{tabular}\end{table}}

\newcommand{\ergebnistabelle}{%
\begin{table}[htbp]\centering\footnotesize\setlength{\tabcolsep}{3.5pt}
\caption{Ergebnisreihe 2017--2025 (Mio.~EUR; EPS und Dividende in EUR; ROACE und
Verschuldungsgrad in \%). Quelle je Jahr ist der j\"ungste Gesch\"aftsbericht, der das Jahr
druckt (Kennzahlenseiten GB~2019 S.~123, GB~2020 S.~129, GB~2021 S.~162, GB~2022 S.~172,
GB~2023 S.~171, GB~2024 S.~150, GB~2025 S.~146). Vergleichbares EBITDA in heutiger
Abgrenzung erst ab 2019. ROACE bis 2019 alte Abgrenzung, ab 2020 \enquote{Comparable ROACE}.}
\label{tab:ergebnis}
\begin{tabular}{@{}lrrrrrrrrr@{}}\toprule
Jahr & Umsatz & Vergl.\ EBITDA & EBIT & Ergebnis & EPS & Dividende & ROACE & Nettoschulden & Verschuldung\\\midrule
\tabellenkoerper{data/reihe_ergebnis}
\bottomrule\end{tabular}\end{table}}

\newcommand{\zahlungstabelle}{%
\begin{table}[htbp]\centering\footnotesize\setlength{\tabcolsep}{3.5pt}
\caption{Zahlungsreihe 2017--2025 und letzte zw\"olf Monate (Mio.~EUR). Freier Zufluss =
operativer Mittelzufluss + Investitionen (Sach- und immaterielle Anlagen) + Leasingtilgung.
Instandhaltung: Kassenabfluss f\"ur Instandhaltungsinvestitionen laut Lagebericht.
Zufluss ohne Wachstum = operativer Mittelzufluss $-$ Instandhaltung + Leasingtilgung.
Leasing 2019 und 2020 aus dem Anhang (GB~2020 S.~193), vor 2019 keine aktivierten
Leasingverh\"altnisse (IAS~17).}\label{tab:zahlung}
\begin{tabular}{@{}lrrrrrrrl@{}}\toprule
Jahr & Oper.\ CF & Investitionen & Leasing & Freier Zufluss & Instandhaltung & ohne Wachstum & Dividende & Beleg oper.\ CF\\\midrule
\tabellenkoerper{data/reihe_zahlung}
\bottomrule\end{tabular}\end{table}}

\newcommand{\prognosetabelle}{%
\begin{table}[htbp]\centering\small
\caption{Prognosetreue der Investitionen ohne M\&A 2021--2025. Gelesen ist die Prognose
\emph{wie zuerst gegeben} (Ausblick im Gesch\"aftsbericht des Vorjahres). \enquote{ca.} wird
als $\pm$\,5\,\% gelesen. Ist-Werte aus dem Lagebericht des jeweiligen Jahres.}\label{tab:prognose}
\begin{tabular}{@{}llrrll@{}}\toprule
Jahr & Prognose (zuerst) & Ist (Mio.~EUR) & Abweichung zur Mitte & Lage & Prognose aus\\\midrule
\tabellenkoerper{data/prognose}
\bottomrule\end{tabular}\end{table}}

\newcommand{\margentabelle}{%
\begin{table}[htbp]\centering\small
\caption{Vergleichbare Verkaufsmarge von Renewable Products (USD/t), vergleichbares EBITDA
des Konzerns (Mio.~EUR) und Schlusskurs (EUR). Marge je Jahr in der j\"ungsten Darstellung
(Einzelbelege im Text); 2022 nach der 2023 ge\"anderten Formel. Kurs zum Jahresende laut
Kennzahlenseite, letzte Zeile Kurs vom \KursDatum{} (Yahoo Finance) gegen Jahresende 2025.}\label{tab:marge}
\begin{tabular}{@{}lrrrrr@{}}\toprule
Jahr & RP-Marge & Ver\"and.\ in \% & Vergl.\ EBITDA & Kurs & Ver\"and.\ in \%\\\midrule
\tabellenkoerper{data/marge}
\bottomrule\end{tabular}\end{table}}

\newcommand{\optionentabelle}{%
\begin{table}[htbp]\centering\small
\caption{Interner Zinsfu{\ss} auf zehn Jahre je Handlungsoption und Zuflussszenario. Option~C
legt denselben Betrag zur H\"urde an; ihre Zeile ist die Messlatte.}\label{tab:optionen}
\begin{tabular}{@{}lrrrrrr@{}}\toprule
Option & Kurs (EUR) & pess. & Basis & Ernte & LTM & LTM o.\,UV\\\midrule
\tabellenkoerper{data/optionen}
\bottomrule\end{tabular}\end{table}}

\newcommand{\schwellentabelle}{%
\begin{table}[htbp]\centering\small
\caption{Schwellenkurs je Aktie (EUR): Kurs, zu dem der interne Zinsfu{\ss} die H\"urde von
\Pz{\Huerde} erreicht; letzte Spalte bei \Pz{\HuerdeLang}. Zufluss konstant ab 2027,
Endwert = Buchwert des Eigenkapitals zum 30.06.2026.}\label{tab:schwelle}
\begin{tabular}{@{}lrrrrr@{}}\toprule
Szenario & Zufluss (Mio.~EUR) & 5 Jahre & 10 Jahre & 15 Jahre & 10 J., \Pz{\HuerdeLang}\\\midrule
\tabellenkoerper{data/schwelle}
\bottomrule\end{tabular}\end{table}}

\newcommand{\umkehrtabellen}{%
\begin{table}[htbp]\centering\small
\caption{Umkehrungen zum heutigen Kurs bei \Pz{\Huerde}: erforderliches Wachstum des freien
Zuflusses p.\,a.\ (oben) und erforderlicher Endwert als Vielfaches des Buchwerts
(unten).}\label{tab:umkehr}
\begin{tabular}{@{}lrrr@{}}\toprule
Szenario & 5 Jahre & 10 Jahre & 15 Jahre\\\midrule
\tabellenkoerper{data/wachstum}
\midrule
\tabellenkoerper{data/endwert}
\bottomrule\end{tabular}\end{table}}

\newcommand{\dauertabelle}{%
\begin{table}[htbp]\centering\small
\caption{Wie lange m\"usste die Sondermarge anhalten? Interner Zinsfu{\ss} auf zehn Jahre zum
heutigen Kurs, wenn der Zufluss $k$ Jahre auf dem LTM-Niveau (\Mio{\FcfLtm}{EUR}) liegt und danach
auf den Median (\Mio{\FcfMedian}{EUR}) zur\"uckf\"allt. Endwert links Buchwert, rechts der heutige
B\"orsenwert (Ausstieg ohne Bewertungs\"anderung).}\label{tab:dauer}
\begin{tabular}{@{}rrr@{}}\toprule
$k$ (Jahre) & Endwert Buchwert & Endwert heutiger B\"orsenwert\\\midrule
\tabellenkoerper{data/dauer}
\bottomrule\end{tabular}\end{table}}

\newcommand{\evtabelle}{%
\begin{table}[htbp]\centering\small
\caption{Unternehmenswert zu vergleichbarem EBITDA: B\"orsenwert und Nettofinanzschulden zum
Jahresende laut Kennzahlenseite, vergleichbares EBITDA des Jahres. Letzte Zeile: B\"orsenwert
zum \KursDatum, Nettofinanzschulden zum 30.06.2026, EBITDA der letzten zw\"olf Monate.}\label{tab:ev}
\begin{tabular}{@{}lrrrr@{}}\toprule
Jahr & B\"orsenwert & Nettoschulden & Vergl.\ EBITDA & EV/EBITDA\\\midrule
\tabellenkoerper{data/evebitda}
\bottomrule\end{tabular}\end{table}}

\newcommand{\kurzfristtabelle}{%
\begin{table}[htbp]\centering\small
\caption{Kurzer Horizont: Kurs je Aktie, wenn der Markt Neste zum Median-Multiple
2019--2025 (\EvMedian{}$\times$) auf dem jeweiligen EBITDA-Niveau bewertet;
Nettofinanzschulden zum 30.06.2026. Abstand zum Kurs vom \KursDatum.}\label{tab:kurzfrist}
\begin{tabular}{@{}lrrr@{}}\toprule
EBITDA-Niveau & Vergl.\ EBITDA (Mio.~EUR) & Kurs (EUR) & Abstand\\\midrule
\tabellenkoerper{data/kurzfrist}
\bottomrule\end{tabular}\end{table}}

\newcommand{\kursabbildung}{%
\begin{figure}[htbp]\centering
\begin{tikzpicture}
\begin{axis}[width=\linewidth, height=6.2cm, date coordinates in=x,
  xticklabel=\year, xtick={2017-01-01,2019-01-01,2021-01-01,2023-01-01,2025-01-01},
  ylabel={EUR}, ymin=0, grid=major, grid style={gray!20},
  tick label style={font=\footnotesize}, label style={font=\footnotesize}]
\addplot[kursfarbe, thick, mark=none] table[x=datum, y=kurs, col sep=comma] {data/kursreihe.csv};
\end{axis}
\end{tikzpicture}
\caption{Neste, Schlusskurs zum Monatsende, Dezember 2016 bis September 2026
(Yahoo Finance, splitbereinigt).}\label{fig:kurs}
\end{figure}}
